\renewcommand{\thesection}{\arabic{section}}\renewcommand{\theHsection}{chapter.\arabic{section}}\setcounter{section}{4}\section{求导的逆运算}
\noindent\begin{minipage}{\textwidth}\ProblemHeading{习题}{5.1}
求下列不定积分；出现的参数为常数，在被积函数有定义的区间上讨论：
\SubQ{1}{$\displaystyle\int (2+x^5)^2\,\mathrm dx$。}
\SubQ{2}{$\displaystyle\int \left(\frac{1-x}{x}\right)^2\,\mathrm dx$。}
\SubQ{3}{$\displaystyle\int (1-1/x^2)\sqrt x\,\mathrm dx$。}
\end{minipage}\par\vfill
\noindent\begin{minipage}{\textwidth}\ProblemHeading{习题}{5.2}
求下列不定积分；出现的参数为常数，在被积函数有定义的区间上讨论：
\SubQ{1}{$\displaystyle\int \cosh x\,\mathrm dx$。}
\SubQ{2}{$\displaystyle\int \sinh x\,\mathrm dx$。}
\SubQ{3}{$\displaystyle\int \frac{2^{x+1}-5^{x-1}}{10^x}\,\mathrm dx$。}
\end{minipage}\par\vfill
\noindent\begin{minipage}{\textwidth}\ProblemHeading{习题}{5.3}
求下列不定积分；出现的参数为常数，在被积函数有定义的区间上讨论：
\SubQ{1}{$\displaystyle\int \frac{\mathrm e^{3x}+1}{\mathrm e^x+1}\,\mathrm dx$。}
\SubQ{2}{$\displaystyle\int \sqrt{1-\sin(2x)}\,\mathrm dx$。}
\SubQ{3}{$\displaystyle\int \frac1{(x+a)(x+b)}\,\mathrm dx$。}
\end{minipage}\par\vfill
\newpage
\noindent\begin{minipage}{\textwidth}\ProblemHeading{习题}{5.4}
求下列不定积分；出现的参数为常数，在被积函数有定义的区间上讨论：
\SubQ{1}{$\displaystyle\int \cos^2x\,\mathrm dx$。}
\SubQ{2}{$\displaystyle\int \sin^2x\,\mathrm dx$。}
\SubQ{3}{$\displaystyle\int \frac{x^5}{1+x}\,\mathrm dx$。}
\end{minipage}\par\vfill
\noindent\begin{minipage}{\textwidth}\ProblemHeading{习题}{5.5}
求下列不定积分；出现的参数为常数，在被积函数有定义的区间上讨论：
\SubQ{1}{$\displaystyle\int \frac1{\cos^2x\sin^2x}\,\mathrm dx$。}
\SubQ{2}{$\displaystyle\int \frac{x^4}{1+x^2}\,\mathrm dx$。}
\end{minipage}\par\vfill
\noindent\begin{minipage}{\textwidth}\ProblemHeading{例题}{5.6}
计算 $\int x\mathrm e^x\,\mathrm dx$。
\end{minipage}\par\vfill
\newpage
\noindent\begin{minipage}{\textwidth}\ProblemHeading{例题}{5.7}
计算 $\int x^2\mathrm e^x\,\mathrm dx$。
\end{minipage}\par\vfill
\noindent\begin{minipage}{\textwidth}\ProblemHeading{例题}{5.8}
计算 $\int\ln x\,\mathrm dx$。
\end{minipage}\par\vfill
\noindent\begin{minipage}{\textwidth}\ProblemHeading{例题}{5.9}
计算下列不定积分：
\SubQ{1}{$\int\mathrm e^x\cos x\,\mathrm dx$。}
\SubQ{2}{$\int\mathrm e^x\sin x\,\mathrm dx$。}
\end{minipage}\par\vfill
\newpage
\noindent\begin{minipage}{\textwidth}\ProblemHeading{例题}{5.10}
设 $n\in\mathbb N^*$。计算下列不定积分，可用递推公式表示：
\SubQ{1}{$\int\cos^n x\,\mathrm dx$。}
\SubQ{2}{$\int\sin^n x\,\mathrm dx$。}
\end{minipage}\par\vfill
\noindent\begin{minipage}{\textwidth}\ProblemHeading{例题}{5.11}
求 $\int\sin^3x\cos x\,\mathrm dx$。
\end{minipage}\par\vfill
\noindent\begin{minipage}{\textwidth}\ProblemHeading{例题}{5.12}
求 $\int x\mathrm e^{x^2}\,\mathrm dx$。
\end{minipage}\par\vfill
\newpage
\noindent\begin{minipage}{\textwidth}\ProblemHeading{例题}{5.13}
求 $\int\mathrm dx/(ax+b)$，其中 $a,b$ 为常数且 $a\ne0$。
\end{minipage}\par\vfill
\noindent\begin{minipage}{\textwidth}\ProblemHeading{例题}{5.14}
求 $\int(\ln x)/x\,\mathrm dx$。
\end{minipage}\par\vfill
\noindent\begin{minipage}{\textwidth}\ProblemHeading{例题}{5.15}
求 $\int\mathrm dx/(a^2+x^2)$，其中常数 $a\ne0$。
\end{minipage}\par\vfill
\newpage
\noindent\begin{minipage}{\textwidth}\ProblemHeading{例题}{5.16}
计算 $\int\dfrac{\mathrm dx}{a^2\sin^2x+b^2\cos^2x}$，其中 $a,b$ 为常数且 $ab\ne0$。
\end{minipage}\par\vfill
\noindent\begin{minipage}{\textwidth}\ProblemHeading{例题}{5.17}
求 $\int\sqrt{a^2-x^2}\,\mathrm dx$，其中 $a>0$、$|x|\le a$。
\end{minipage}\par\vfill
\noindent\begin{minipage}{\textwidth}\ProblemHeading{例题}{5.18}
求 $\int\mathrm dx/\sqrt{a^2+x^2}$，其中常数 $a>0$。
\end{minipage}\par\vfill
\newpage
\noindent\begin{minipage}{\textwidth}\ProblemHeading{例题}{5.19}
求 $\int\sqrt{a^2+x^2}\,\mathrm dx$，其中 $a$ 为常数。
\end{minipage}\par\vfill
\noindent\begin{minipage}{\textwidth}\ProblemHeading{例题}{5.20}
求 $\int\mathrm dx/(1+\sqrt{1+x})$。
\end{minipage}\par\vfill
\noindent\begin{minipage}{\textwidth}\ProblemHeading{例题}{5.21}
求 $\int\dfrac{\sin\sqrt{x}}{\sqrt{x}}\,\mathrm dx$。
\end{minipage}\par\vfill
\newpage
\noindent\begin{minipage}{\textwidth}\ProblemHeading{习题}{5.22}
求下列不定积分；出现的参数为常数，在被积函数有定义的区间上讨论：
\SubQ{1}{$\displaystyle\int \arctan x\,\mathrm dx$。}
\SubQ{2}{$\displaystyle\int \arcsin x\,\mathrm dx$。}
\SubQ{3}{$\displaystyle\int x\operatorname{arccot}x\,\mathrm dx$。}
\end{minipage}\par\vfill
\noindent\begin{minipage}{\textwidth}\ProblemHeading{习题}{5.23}
求下列不定积分；出现的参数为常数，在被积函数有定义的区间上讨论：
\SubQ{1}{$\displaystyle\int x^2\arccos x\,\mathrm dx$。}
\SubQ{2}{$\displaystyle\int \ln(x+\sqrt{1+x^2})\,\mathrm dx$。}
\SubQ{3}{$\displaystyle\int \sqrt{x}\ln^2x\,\mathrm dx$。}
\end{minipage}\par\vfill
\noindent\begin{minipage}{\textwidth}\ProblemHeading{习题}{5.24}
求下列不定积分；出现的参数为常数，在被积函数有定义的区间上讨论：
\SubQ{1}{$\displaystyle\int x^2\cos x\,\mathrm dx$。}
\SubQ{2}{$\displaystyle\int \frac{x}{\cos^2x}\,\mathrm dx$。}
\SubQ{3}{$\displaystyle\int \frac{\arctan x}{x^2}\,\mathrm dx$。}
\end{minipage}\par\vfill
\newpage
\noindent\begin{minipage}{\textwidth}\ProblemHeading{习题}{5.25}
求下列不定积分；出现的参数为常数，在被积函数有定义的区间上讨论：
\SubQ{1}{$\displaystyle\int x^2\cosh x\,\mathrm dx$。}
\SubQ{2}{$\displaystyle\int \cos(\ln x)\,\mathrm dx$。}
\SubQ{3}{$\displaystyle\int \sin(\ln x)\,\mathrm dx$。}
\end{minipage}\par\vfill
\noindent\begin{minipage}{\textwidth}\ProblemHeading{习题}{5.26}
设 $p$ 为 $n$ 次多项式，$a$ 为常数。用 $p$ 的各阶导数表示原函数 $\int p(x)\mathrm e^{ax}\,\mathrm dx$；对 $a=0$ 的情形另作说明。
\end{minipage}\par\vfill
\noindent\begin{minipage}{\textwidth}\ProblemHeading{习题}{5.27}
设 $f$ 二阶可导。计算 $\int xf^{\prime\prime}(x)\,\mathrm dx$。
\end{minipage}\par\vfill
\newpage
\noindent\begin{minipage}{\textwidth}\ProblemHeading{习题}{5.28}
求下列不定积分；出现的参数为常数，在被积函数有定义的区间上讨论：
\SubQ{1}{$ \int x e^{-x^2} dx $;}
\SubQ{2}{$ \int\frac{x}{\sqrt{1-x^{2}}}dx $;}
\SubQ{3}{$ \int \frac{\ln^{2}x}{x} dx $;}
\end{minipage}\par\vfill
\noindent\begin{minipage}{\textwidth}\ProblemHeading{习题}{5.29}
求下列不定积分；出现的参数为常数，在被积函数有定义的区间上讨论：
\SubQ{1}{$ \int \frac{x}{x^{2} + 4} \, dx $;}
\SubQ{2}{$ \int \frac{dx}{\cosh x} $;}
\SubQ{3}{$ \int \frac{dx}{\sqrt{x}(1+x)}; $}
\end{minipage}\par\vfill
\noindent\begin{minipage}{\textwidth}\ProblemHeading{习题}{5.30}
求下列不定积分；出现的参数为常数，在被积函数有定义的区间上讨论：
\SubQ{1}{$ \int \frac{dx}{\sin x} $;}
\SubQ{2}{$ \int \frac{dx}{x \ln x (\ln \ln x)}; $}
\SubQ{3}{$ \int \frac{1}{x^{2}} \sin \frac{1}{x} dx $;}
\end{minipage}\par\vfill
\newpage
\noindent\begin{minipage}{\textwidth}\ProblemHeading{习题}{5.31}
求下列不定积分；出现的参数为常数，在被积函数有定义的区间上讨论：
\SubQ{1}{$ \int \frac{\arctan^{2}x}{1 + x^{2}} \mathrm{d}x $;}
\SubQ{2}{$ \int \frac{dx}{(a + bx)^2} (ab \neq 0) $;}
\SubQ{3}{$ \int \cos ax \sin bx \, dx $;}
\end{minipage}\par\vfill
\noindent\begin{minipage}{\textwidth}\ProblemHeading{习题}{5.32}
求下列不定积分；出现的参数为常数，在被积函数有定义的区间上讨论：
\SubQ{1}{$ \int \cos ax \cos bx \, dx $;}
\SubQ{2}{$\int \sin ax \sin bx \, dx$;}
\SubQ{3}{$ \int \sin^{4} x \, dx $;}
\end{minipage}\par\vfill
\noindent\begin{minipage}{\textwidth}\ProblemHeading{习题}{5.33}
求下列不定积分；出现的参数为常数，在被积函数有定义的区间上讨论：
\SubQ{1}{$ \int \sin^{5} x \, dx $;}
\SubQ{2}{$ \int \frac{\cos x + \sin x}{\sqrt[3]{\sin x - \cos x}} \, dx $;}
\SubQ{3}{$ \int \frac{dx}{2 - \sin^2 x} $;}
\end{minipage}\par\vfill
\newpage
\noindent\begin{minipage}{\textwidth}\ProblemHeading{习题}{5.34}
求下列不定积分；出现的参数为常数，在被积函数有定义的区间上讨论：
\SubQ{1}{$ \int \frac{\cos x \sin x}{a^{2} \cos^{2} x + b^{2} \sin^{2} x} \, dx $;}
\SubQ{2}{$ \int \frac{dx}{\cos x + \sin x} $;}
\SubQ{3}{$ \int \frac{dx}{a\cos x + b\sin x} $;}
\end{minipage}\par\vfill
\noindent\begin{minipage}{\textwidth}\ProblemHeading{习题}{5.35}
求下列不定积分；出现的参数为常数，在被积函数有定义的区间上讨论：
\SubQ{1}{$ \int \frac{\cos x}{\sqrt{2 + \cos 2x}} \, dx $;}
\SubQ{2}{$ \int \frac{\ln x}{x \sqrt{1 + \ln x}} \, dx $;}
\SubQ{3}{$ \int \frac{x^{3}}{\sqrt{1-x^{2}}} dx $;}
\end{minipage}\par\vfill
\noindent\begin{minipage}{\textwidth}\ProblemHeading{习题}{5.36}
求下列不定积分；出现的参数为常数，在被积函数有定义的区间上讨论：
\SubQ{1}{$ \int \frac{x \, dx}{\sqrt{1 + x^{2} + (1 + x^{2})^{3/2}}} $;}
\SubQ{2}{$ \int \frac{dx}{x \sqrt{x^2 + 1}}; $}
\SubQ{3}{$ \int e^{\sqrt{x}} \, dx $;}
\end{minipage}\par\vfill
\newpage
\noindent\begin{minipage}{\textwidth}\ProblemHeading{习题}{5.37}
求下列不定积分；出现的参数为常数，在被积函数有定义的区间上讨论：
\SubQ{1}{$ \int x^{3} e^{-x^{2}} \, dx $;}
\SubQ{2}{$ \int \arctan \sqrt{x} \, dx $;}
\SubQ{3}{$ \int x \sin \sqrt{x} \, dx $;}
\end{minipage}\par\vfill
\noindent\begin{minipage}{\textwidth}\ProblemHeading{习题}{5.38}
求下列不定积分；出现的参数为常数，在被积函数有定义的区间上讨论：
\SubQ{1}{$ \int \frac{x^{n/2}}{\sqrt{1 + x^{n+2}}} \, dx \, (n \in \mathbb{N}^*) \, ; $}
\SubQ{2}{$ \int \frac{x^2}{\sqrt{x^2 + a^2}} \, dx \, ; $}
\SubQ{3}{$ \int \frac{\cos x \sin x}{1 + \sin^{4} x} \, dx $.}
\end{minipage}\par\vfill
\noindent\begin{minipage}{\textwidth}\ProblemHeading{习题}{5.39}
设 $I_n=\int\ln^nx\,\mathrm dx$，$n\in\mathbb N^*$。推导 $I_n$ 的递推公式。
\end{minipage}\par\vfill
\newpage
\noindent\begin{minipage}{\textwidth}\ProblemHeading{习题}{5.40}
取常数 $a>0$，用换元 $x=a\sinh t$ 或 $x=a\cosh t$，在相应代换覆盖的区间内计算：
\SubQ{1}{$\int(x^2+a^2)^{-3/2}\,\mathrm dx$。}
\SubQ{2}{$\int\dfrac{x^2}{\sqrt{x^2-a^2}}\,\mathrm dx$。}
\SubQ{3}{$\int\sqrt{x^2-a^2}\,\mathrm dx$。}
\end{minipage}\par\vfill
\noindent\begin{minipage}{\textwidth}\ProblemHeading{习题}{5.41}
求下列不定积分；出现的参数为常数，在被积函数有定义的区间上讨论：
\SubQ{1}{$\displaystyle\int (1-2/x)^2\mathrm e^x\,\mathrm dx$。}
\SubQ{2}{$\displaystyle\int \frac{x\mathrm e^x}{(1+x)^2}\,\mathrm dx$。}
\end{minipage}\par\vfill
\noindent\begin{minipage}{\textwidth}\ProblemHeading{习题}{5.42}
分别求出在所给条件涉及的自变量范围内的函数 $f$：
\SubQ{1}{$f^{\prime}(x^2)=1/x$（$x>0$）。}
\SubQ{2}{$f^{\prime}(\sin^2x)=\cos^2x$（$x\in\mathbb R$）。}
\end{minipage}\par\vfill
\newpage
\noindent\begin{minipage}{\textwidth}\ProblemHeading{例题}{5.43}
将 $\dfrac{x+1}{x^2-4x+3}$ 分解成部分分式。
\end{minipage}\par\vfill
\noindent\begin{minipage}{\textwidth}\ProblemHeading{例题}{5.44}
将 $\dfrac{x}{x^3+x^2+3x+3}$ 分解成部分分式。
\end{minipage}\par\vfill
\noindent\begin{minipage}{\textwidth}\ProblemHeading{例题}{5.45}
将 $\dfrac{x^3+1}{x^4-3x^3+3x^2-x}$ 分解成部分分式。
\end{minipage}\par\vfill
\newpage
\noindent\begin{minipage}{\textwidth}\ProblemHeading{例题}{5.46}
计算 $\int\dfrac{x^3+1}{x^4-3x^3+3x^2-x}\,\mathrm dx$。
\end{minipage}\par\vfill
\noindent\begin{minipage}{\textwidth}\ProblemHeading{例题}{5.47}
计算 $\int\dfrac{5x+6}{x^2+x+1}\,\mathrm dx$。
\end{minipage}\par\vfill
\noindent\begin{minipage}{\textwidth}\ProblemHeading{例题}{5.48}
计算 $\int\dfrac{5x+3}{(x^2-2x+5)^2}\,\mathrm dx$。
\end{minipage}\par\vfill
\newpage
\noindent\begin{minipage}{\textwidth}\ProblemHeading{习题}{5.49}
求下列不定积分；出现的参数为常数，在被积函数有定义的区间上讨论：
\SubQ{1}{$\displaystyle\int \frac{x}{2x^2-3x-2}\,\mathrm dx$。}
\SubQ{2}{$\displaystyle\int \frac{2x^2+1}{(x+3)(x-1)(x-4)}\,\mathrm dx$。}
\SubQ{3}{$\displaystyle\int \frac1{2x^3+3x^2+x}\,\mathrm dx$。}
\end{minipage}\par\vfill
\noindent\begin{minipage}{\textwidth}\ProblemHeading{习题}{5.50}
求下列不定积分；出现的参数为常数，在被积函数有定义的区间上讨论：
\SubQ{1}{$\displaystyle\int \frac{x^2}{x^3+5x^2+8x+4}\,\mathrm dx$。}
\SubQ{2}{$\displaystyle\int \frac{4x+3}{(x-2)^3}\,\mathrm dx$。}
\SubQ{3}{$\displaystyle\int \frac{x^2}{(x+2)^2(x+4)^2}\,\mathrm dx$。}
\end{minipage}\par\vfill
\noindent\begin{minipage}{\textwidth}\ProblemHeading{习题}{5.51}
求下列不定积分；出现的参数为常数，在被积函数有定义的区间上讨论：
\SubQ{1}{$\displaystyle\int \frac{8x^3+7}{(x+1)(2x+1)^3}\,\mathrm dx$。}
\SubQ{2}{$\displaystyle\int \frac1{(x+1)(x+2)^2(x+3)^3}\,\mathrm dx$。}
\SubQ{3}{$\displaystyle\int \frac1{1+x^3}\,\mathrm dx$。}
\end{minipage}\par\vfill
\newpage
\noindent\begin{minipage}{\textwidth}\ProblemHeading{习题}{5.52}
求下列不定积分；出现的参数为常数，在被积函数有定义的区间上讨论：
\SubQ{1}{$\displaystyle\int \frac1{(x^2+9)^3}\,\mathrm dx$。}
\SubQ{2}{$\displaystyle\int \frac1{x(x^2+1)^2}\,\mathrm dx$。}
\SubQ{3}{$\displaystyle\int \frac1{1+x^4}\,\mathrm dx$。}
\end{minipage}\par\vfill
\noindent\begin{minipage}{\textwidth}\ProblemHeading{习题}{5.53}
求下列不定积分；出现的参数为常数，在被积函数有定义的区间上讨论：
\SubQ{1}{$\displaystyle\int \frac{1+x^2}{1+x^4}\,\mathrm dx$。}
\SubQ{2}{$\displaystyle\int \frac1{x^4(1+x^2)}\,\mathrm dx$。}
\end{minipage}\par\vfill
\noindent\begin{minipage}{\textwidth}\ProblemHeading{例题}{5.54}
求 $\int\dfrac{\mathrm dx}{\sin x(1+\cos x)}$。
\end{minipage}\par\vfill
\newpage
\noindent\begin{minipage}{\textwidth}\ProblemHeading{例题}{5.55}
计算 $\int\dfrac{\sin^4x}{\cos^2x}\,\mathrm dx$。
\end{minipage}\par\vfill
\noindent\begin{minipage}{\textwidth}\ProblemHeading{例题}{5.56}
设 $a<b$ 为常数，计算 $\int\sqrt{\dfrac{x-a}{b-x}}\,\mathrm dx$，其中 $a<x<b$。
\end{minipage}\par\vfill
\noindent\begin{minipage}{\textwidth}\ProblemHeading{例题}{5.57}
计算 $\int\dfrac{\mathrm dx}{\sqrt{x}(1+\sqrt[3]{x})}$。
\end{minipage}\par\vfill
\newpage
\noindent\begin{minipage}{\textwidth}\ProblemHeading{例题}{5.58}
计算 $\int\dfrac{\sqrt[3]{1+\sqrt[4]{x}}}{\sqrt{x}}\,\mathrm dx$。
\end{minipage}\par\vfill
\noindent\begin{minipage}{\textwidth}\ProblemHeading{例题}{5.59}
计算 $\int\dfrac{\mathrm dx}{\sqrt[4]{1+x^4}}$。
\end{minipage}\par\vfill
\noindent\begin{minipage}{\textwidth}\ProblemHeading{习题}{5.60}
求下列不定积分，在被积函数有定义的区间上讨论：
\SubQ{1}{$ \int \frac{\sin^{3}x}{\cos^{4}x} dx $;}
\SubQ{2}{$ \int \frac{dx}{\cos^4 x \sin x} $;}
\SubQ{3}{$ \int \frac{dx}{2\sin x - \cos x + 5} $;}
\end{minipage}\par\vfill
\newpage
\noindent\begin{minipage}{\textwidth}\ProblemHeading{习题}{5.61}
求下列不定积分，在被积函数有定义的区间上讨论：
\SubQ{1}{$ \int \frac{\cos x \sin x}{\cos x + \sin x} \, dx $;}
\SubQ{2}{$ \int \frac{\sin x}{1 + \cos x + \sin x} \, dx $;}
\SubQ{3}{$ \int \frac{\sin^{2}x}{1 + \sin^{2}x} \, dx $;}
\end{minipage}\par\vfill
\noindent\begin{minipage}{\textwidth}\ProblemHeading{习题}{5.62}
求下列不定积分，在被积函数有定义的区间上讨论：
\SubQ{1}{$ \int \frac{dx}{\cos^4 x + \sin^4 x} $;}
\SubQ{2}{$ \int \frac{dx}{\cos^4 x \sin^4 x} $;}
\SubQ{3}{$ \int \frac{\sin x}{\cos^3 x + \sin^3 x} \, dx $;}
\end{minipage}\par\vfill
\noindent\begin{minipage}{\textwidth}\ProblemHeading{习题}{5.63}
求下列不定积分，在被积函数有定义的区间上讨论：
\SubQ{1}{$ \int \sqrt{\tan x} \, dx $;}
\SubQ{2}{$ \int \frac{dx}{1 + \varepsilon \cos x} $ (常数  $ \varepsilon > 1 $).}
\end{minipage}\par\vfill
\newpage
\noindent\begin{minipage}{\textwidth}\ProblemHeading{习题}{5.64}
求下列不定积分；出现的参数为常数，在被积函数有定义的区间上讨论：
\SubQ{1}{$ \int \frac{dx}{1 + \sqrt[3]{1 + x}}; $}
\SubQ{2}{$ \int \frac{\sqrt[3]{x}}{x(\sqrt{x} + \sqrt[3]{x})} \, dx $;}
\SubQ{3}{$ \int\frac{x}{\sqrt{x+1}+\sqrt[3]{x+1}}dx $;}
\end{minipage}\par\vfill
\noindent\begin{minipage}{\textwidth}\ProblemHeading{习题}{5.65}
求下列不定积分；出现的参数为常数，在被积函数有定义的区间上讨论：
\SubQ{1}{$ \int \frac{\sqrt{1 + x^2}}{2 + x^2} \, dx $;}
\SubQ{2}{$ \int\frac{dx}{x\sqrt{1+x^{2}}} $;}
\SubQ{3}{$ \int \frac{x^3}{\sqrt{1 + 2x^2}} \, dx $;}
\end{minipage}\par\vfill
\noindent\begin{minipage}{\textwidth}\ProblemHeading{习题}{5.66}
求下列不定积分；出现的参数为常数，在被积函数有定义的区间上讨论：
\SubQ{1}{$ \int \sqrt{\frac{1 - \sqrt{x}}{1 + \sqrt{x}}} \, dx $;}
\SubQ{2}{$ \int \frac{dx}{x \sqrt{a^2 - x^2}} (|x| < a, x \neq 0) $;}
\SubQ{3}{$ \int \sqrt{\frac{1 - x}{1 + x}} \frac{dx}{x^2} $;}
\end{minipage}\par\vfill
\newpage
\noindent\begin{minipage}{\textwidth}\ProblemHeading{习题}{5.67}
求下列不定积分；出现的参数为常数，在被积函数有定义的区间上讨论：
\SubQ{1}{$ \int \frac{dx}{(x + a)^2(x + b)^3} $;}
\SubQ{2}{$\displaystyle\int\frac{\mathrm dx}{(x+a)^m(x+b)^n}$，$m,n\in\mathbb N^*$。}
\SubQ{3}{$\displaystyle\int\frac{\mathrm dx}{(1+x^n)\sqrt[n]{1+x^n}}$，$n\in\mathbb N^*$。}
\end{minipage}\par\vfill
